% !TeX root = ../../main.tex

\section{Casos de uso de Big Data}

\subsection{Introducción}

Big Data permite a las organizaciones recopilar, almacenar y analizar grandes cantidades de información. Estos datos pueden proceder de ventas, redes sociales, aplicaciones, sensores, búsquedas y actividades realizadas por los usuarios.

El objetivo no es solamente guardar información, sino utilizarla para detectar problemas, conocer el comportamiento de los clientes, hacer predicciones y tomar decisiones. Los siguientes casos muestran cómo Walmart, Netflix, Microsoft, Facebook, LinkedIn y Airbnb utilizan Big Data en diferentes actividades.

Las cantidades y resultados mencionados corresponden al periodo analizado en los casos de estudio.

\subsection{Walmart}

\textbf{Problema que busca resolver}

Walmart tiene miles de tiendas y vende millones de productos. Uno de sus principales problemas es conseguir que cada establecimiento tenga los productos correctos, en la cantidad adecuada y en el momento en que los clientes los necesitan.

También debe detectar rápidamente errores de precios, productos agotados, problemas de distribución y cambios inesperados en la demanda.

\textbf{Datos utilizados}

Walmart analiza datos de ventas, inventarios, precios, clima, economía, redes sociales, costos de combustible y eventos cercanos a sus tiendas. También estudia las búsquedas realizadas por los clientes en su página web.

La empresa puede relacionar estos datos para descubrir comportamientos que no serían evidentes si cada fuente fuera analizada por separado. Por ejemplo, las condiciones climáticas pueden provocar un aumento en la venta de determinados alimentos o productos de emergencia.

\textbf{Aplicación de Big Data}

Walmart creó el Data Café, un centro de análisis que recibe información procedente de diferentes áreas de la empresa. Este sistema permite monitorear las ventas y otros indicadores casi en tiempo real.

Cuando el sistema detecta un resultado fuera de lo normal, puede generar una alerta para que los empleados investiguen la causa. De esta manera, la empresa puede corregir un precio equivocado, reponer mercancía o colocar en los estantes un producto que no estaba disponible para los clientes.

\textbf{Tecnologías utilizadas}

Entre las tecnologías mencionadas en el caso se encuentran Hadoop, Spark, Cassandra, R y SAS. Estas herramientas permiten almacenar y analizar grandes cantidades de datos provenientes de las tiendas y de fuentes externas.

\textbf{Resultados y análisis}

El Data Café redujo el tiempo necesario para detectar un problema y proponer una solución. Un proceso que anteriormente podía tardar entre dos y tres semanas pasó a realizarse en aproximadamente 20 minutos.

Este caso demuestra que el análisis en tiempo real puede mejorar el inventario, los precios y la disponibilidad de productos. Sus principales retos son administrar el crecimiento de los datos y contar con personal capacitado.

\subsection{Netflix}

\textbf{Problema que busca resolver}

Netflix necesita ofrecer contenido que sea interesante para cada usuario. Si una persona no encuentra rápidamente algo que desea ver, puede considerar que el servicio no tiene suficiente valor y cancelar su suscripción.

Por esta razón, Netflix intenta predecir qué películas o series pueden gustarle a cada usuario y busca mejorar constantemente la calidad de la reproducción.

\textbf{Datos utilizados}

Netflix analiza los títulos vistos, las búsquedas, las calificaciones, el horario de reproducción y el tiempo utilizado para seleccionar una película. También registra cuándo una persona pausa, adelanta, abandona o termina un contenido.

Para evaluar la calidad del servicio, analiza la ubicación del usuario, la velocidad de transmisión, la calidad de imagen y los problemas de almacenamiento temporal o \textit{buffering}.

\textbf{Aplicación de Big Data}

La empresa utiliza algoritmos de recomendación para comparar el comportamiento de cada persona con el de otros usuarios. Con esta información crea una pantalla personalizada y sugiere contenido con mayores probabilidades de ser reproducido.

Netflix también clasifica las películas y series mediante una gran cantidad de características. Esto le permite crear categorías muy específicas y comprender mejor los gustos de sus usuarios.

Los datos también apoyan las decisiones relacionadas con la producción de contenido original. La empresa puede analizar qué actores, directores, géneros y temas tienen mayor aceptación antes de invertir en una producción.

\textbf{Tecnologías utilizadas}

El caso menciona tecnologías como AWS, Hadoop, Hive, Pig, Cassandra, NoSQL, Teradata y herramientas de aprendizaje automático. Su contenido se distribuye desde diferentes ubicaciones para reducir retrasos durante la reproducción.

\textbf{Resultados y análisis}

Las recomendaciones personalizadas contribuyen a que los usuarios encuentren contenido y permanezcan más tiempo en la plataforma. El análisis de datos también ayudó a la empresa a tomar decisiones sobre producciones originales.

Uno de los principales retos fue analizar información no estructurada, como imágenes, audio y video. Para resolverlo, Netflix clasificó el contenido mediante etiquetas y posteriormente comenzó a utilizar reconocimiento de imágenes y otras técnicas automáticas.

\subsection{Microsoft}

\textbf{Problema que busca resolver}

Muchas empresas desean utilizar Big Data, pero no tienen servidores, programas o especialistas suficientes para construir una plataforma propia. Microsoft busca facilitar el acceso a estas tecnologías mediante servicios en la nube y herramientas de análisis.

\textbf{Datos utilizados}

Microsoft recopila información relacionada con el uso de sus programas y servicios. Esto puede incluir errores del sistema, tiempos de respuesta, funciones utilizadas, características del equipo y datos de interacción con sus aplicaciones.

También puede utilizar información de navegación y preferencias para personalizar servicios y ofrecer publicidad.

\textbf{Aplicación de Big Data}

Microsoft ofrece almacenamiento y procesamiento mediante Azure. Las organizaciones pueden guardar información, ejecutar modelos analíticos y utilizar aprendizaje automático sin adquirir toda la infraestructura física.

La empresa también integró herramientas de análisis en productos conocidos como Excel y Power BI. Esto permite que personas sin conocimientos avanzados de programación creen informes y visualizaciones.

Los datos de funcionamiento de Windows y otros programas son utilizados para detectar fallas, mejorar el rendimiento y conocer qué funciones resultan más útiles para los usuarios.

\textbf{Tecnologías utilizadas}

Entre las tecnologías mencionadas se encuentran Microsoft Azure, SQL Server, HDInsight, Hadoop, Power BI y Excel. Azure también permite desarrollar proyectos relacionados con el Internet de las cosas.

\textbf{Resultados y análisis}

Microsoft convirtió los servicios de nube, almacenamiento y análisis en una parte importante de su negocio. Su principal ventaja es que puede acercar Big Data a empresas pequeñas y medianas mediante herramientas relativamente conocidas.

Uno de sus principales retos es la privacidad. La empresa debe explicar claramente qué información recopila, cómo la utiliza y cuáles datos se destinan a mejorar los programas o a personalizar publicidad.

\subsection{Facebook}

\textbf{Problema que busca resolver}

Las empresas necesitan mostrar sus anuncios a personas que realmente puedan estar interesadas en sus productos. La publicidad tradicional puede ser costosa porque se presenta a una audiencia muy amplia y no siempre llega a los clientes adecuados.

Facebook utiliza Big Data para conocer los intereses generales de sus usuarios y ofrecer publicidad dirigida a grupos específicos.

\textbf{Datos utilizados}

Facebook analiza información demográfica, ubicación, publicaciones, reacciones, páginas seguidas, amistades, fotografías e interacciones. Cada vez que un usuario comparte contenido o selecciona ``Me gusta'', la plataforma obtiene información sobre sus posibles intereses.

La empresa también obtiene datos de las aplicaciones que funcionan dentro de su plataforma y de otros servicios relacionados.

\textbf{Aplicación de Big Data}

Los algoritmos agrupan a los usuarios de acuerdo con características e intereses similares. Después, los anunciantes pueden seleccionar el tipo de público al que desean mostrar sus productos.

Por ejemplo, una empresa puede dirigir un anuncio a personas de cierta ciudad, edad o profesión que también hayan mostrado interés en un tema determinado. Esto permite utilizar el presupuesto de publicidad de manera más eficiente.

\textbf{Tecnologías utilizadas}

El caso menciona el uso de MySQL, Hadoop, HBase y Hive. Facebook también utiliza centros de datos y servidores diseñados para procesar la enorme cantidad de información generada por sus usuarios.

\textbf{Resultados y análisis}

La publicidad segmentada se convirtió en una de las principales fuentes de ingresos de Facebook. También permitió que empresas pequeñas mostraran sus productos a públicos específicos sin contratar campañas masivas.

El reto más importante es proteger la información personal. La plataforma debe evitar el acceso no autorizado, explicar sus políticas y conservar la confianza de los usuarios.

\subsection{LinkedIn}

\textbf{Problema que busca resolver}

LinkedIn debe mantener su plataforma útil para estudiantes, trabajadores, empresas y reclutadores. Para hacerlo, necesita recomendar contactos, empleos, noticias y oportunidades profesionales relacionadas con los intereses de cada usuario.

\textbf{Datos utilizados}

LinkedIn registra clics, visitas a perfiles, mensajes, publicaciones, empleos consultados y contactos agregados. También considera información como empresas anteriores, estudios, ubicación, habilidades y conexiones compartidas.

Con autorización, puede utilizar la lista de contactos del correo electrónico para sugerir personas conocidas.

\textbf{Aplicación de Big Data}

Una de sus aplicaciones más conocidas es la función ``Personas que quizá conozcas''. Los algoritmos analizan contactos compartidos, empresas, escuelas y visitas a perfiles para crear recomendaciones.

LinkedIn también utiliza los datos para recomendar vacantes y mostrar información sobre las personas que visitaron un perfil. En el área comercial, analiza la respuesta de los usuarios a los anuncios para ayudar a los anunciantes a mejorar sus campañas.

La información se procesa continuamente, por lo que los cambios de empleo, las publicaciones y las nuevas conexiones pueden mostrarse rápidamente.

\textbf{Tecnologías utilizadas}

LinkedIn utiliza Hadoop, MapReduce, Kafka, Pig, Hive, Oracle, Java y MySQL. También desarrolló herramientas propias como Voldemort, Espresso y Pinot para almacenar y analizar información.

\textbf{Resultados y análisis}

Las recomendaciones personalizadas contribuyeron al crecimiento de la plataforma y a que los usuarios construyeran redes profesionales más útiles. Big Data también apoya sus servicios de contratación, membresías y publicidad.

Sus retos incluyen procesar una cantidad creciente de información, contratar especialistas y ser transparente sobre el uso de los datos personales.

\subsection{Airbnb}

\textbf{Problema que busca resolver}

Airbnb necesita relacionar a millones de viajeros con anfitriones que ofrecen habitaciones, departamentos o casas. La propiedad recomendada debe ajustarse a la ubicación, precio, fechas y preferencias del huésped.

La empresa también debe ayudar a los anfitriones a establecer un precio adecuado y detectar reservaciones o pagos que puedan ser fraudulentos.

\textbf{Datos utilizados}

Airbnb analiza búsquedas, reservaciones, ubicación, época del año, fotografías, precio, cantidad de habitaciones, servicios disponibles, opiniones y calificaciones.

También puede considerar acontecimientos externos. Por ejemplo, un festival o evento importante puede aumentar la demanda y modificar el precio de los alojamientos en una ciudad.

\textbf{Aplicación de Big Data}

Airbnb utiliza algoritmos para recomendar alojamientos y calcular precios. Su plataforma Aerosolve analiza variables como ubicación, temporada, características de la vivienda, fotografías y cantidad de opiniones.

La empresa también desarrolló Airpal, una herramienta que permite a sus empleados consultar información sin depender completamente del equipo de científicos de datos.

Los algoritmos de aprendizaje automático ayudan a identificar operaciones anormales antes de procesar el pago. El sistema de calificaciones también permite aumentar la confianza entre huéspedes y anfitriones.

\textbf{Tecnologías utilizadas}

El caso menciona Hadoop, HDFS, Hive, Amazon EC2, Presto, Airpal y Aerosolve. Estas tecnologías permiten procesar datos estructurados, como reservaciones, y datos no estructurados, como fotografías y comentarios.

\textbf{Resultados y análisis}

El uso de datos mejoró la asignación de precios, las recomendaciones y la detección de fraudes. También permitió que empleados de diferentes departamentos utilizaran información para tomar decisiones.

El principal reto fue adaptar la infraestructura al rápido crecimiento de la empresa. Airbnb tuvo que desarrollar herramientas sencillas para que los empleados pudieran consultar datos y dejar que los especialistas se concentraran en problemas más importantes.

\subsection{Análisis general}

Aunque las seis empresas pertenecen a sectores diferentes, siguen un proceso parecido:

\begin{enumerate}
    \item Recopilan información de usuarios, operaciones o fuentes externas.
    \item Almacenan los datos en sistemas distribuidos o servicios en la nube.
    \item Utilizan algoritmos para encontrar patrones y hacer predicciones.
    \item Presentan los resultados mediante recomendaciones, alertas o informes.
    \item Utilizan esos resultados para mejorar productos, servicios y decisiones.
\end{enumerate}

Walmart utiliza Big Data para mejorar sus tiendas; Netflix para recomendar y producir contenido; Microsoft para ofrecer herramientas de análisis; Facebook para dirigir publicidad; LinkedIn para conectar profesionales; y Airbnb para relacionar huéspedes con alojamientos.

\subsection{Conclusión}

Estos casos demuestran que Big Data puede aplicarse en ventas, entretenimiento, tecnología, publicidad, empleo y hospedaje. En todos los casos, la información se convierte en acciones que mejoran la eficiencia y la experiencia del cliente.

Sin embargo, también existen retos. Las empresas deben proteger la privacidad, mantener datos confiables, contratar personal capacitado y utilizar infraestructura capaz de crecer. El éxito de Big Data depende tanto de la tecnología como de la capacidad de convertir los resultados en decisiones útiles.